\documentclass[11pt]{amsart}

% Encoding and fonts
\usepackage[T1]{fontenc}
\usepackage[utf8]{inputenc}
\usepackage{lmodern}

% Math packages
\usepackage{amsmath}
\usepackage{amssymb}
\usepackage{amsfonts}
\usepackage{amsthm}
\usepackage{mathtools}

% Layout and utilities
\usepackage{geometry}
\usepackage{enumitem}
\usepackage[hidelinks]{hyperref}

\geometry{margin=1in}

\numberwithin{equation}{section}
\setcounter{tocdepth}{3}
\setcounter{secnumdepth}{3}

% ---------------------------------------------------------
% Theorem-like environments
% ---------------------------------------------------------

\theoremstyle{plain}
\newtheorem{theorem}{Theorem}[section]
\newtheorem{proposition}[theorem]{Proposition}
\newtheorem{lemma}[theorem]{Lemma}
\newtheorem{corollary}[theorem]{Corollary}
\newtheorem{conjecture}[theorem]{Conjecture}

\theoremstyle{definition}
\newtheorem{definition}[theorem]{Definition}
\newtheorem{experiment}[theorem]{Experiment}
\newtheorem{heuristic}[theorem]{Heuristic}
\newtheorem{question}[theorem]{Question}

\theoremstyle{remark}
\newtheorem{remark}[theorem]{Remark}
\newtheorem{warning}[theorem]{Warning}

\DeclareMathOperator{\tr}{tr}
\DeclareMathOperator{\rank}{rank}
\DeclareMathOperator{\Var}{Var}
\DeclareMathOperator{\Cov}{Cov}

\title{
Orientation-Bit Geometry of the M\"obius Path:
Reflection Fibres, a Divisor-Defect Gram Form,
and Exploratory Routes toward Sine-Kernel Statistics
and Two-Point Cancellation
}

\author{}
\date{}

\begin{document}

\maketitle

\tableofcontents

\bigskip

\begin{remark}
\label{rem:exploratory_nature}

This notebook is exploratory in nature.

Its purpose is to develop a single geometric object, the Gram form of
the divisor defect on the orientation bits of the M\"obius path, and to
examine how far that object can be pushed toward sine-kernel level
statistics and toward two-point cancellation for the M\"obius function.

Several sections contain formal manipulations, heuristic arguments and
speculative interpretations.

Whenever a statement is rigorous it is stated inside a theorem, lemma,
proposition or corollary environment and a proof is supplied.

All remaining claims should be regarded as tentative, heuristic or
conjectural, and are marked as such.

No numerical computation is performed anywhere in this notebook. The
few empirical quantities that appear are recorded in
Remark~\ref{rem:empirical_inputs} and are inputs from separate work,
not results obtained here.

\end{remark}

\begin{warning}
\label{warn:bibliographic_metadata}

Bibliographic metadata in Section~\ref{sec:bibliography} has not been
verified against primary sources in this version. DOIs are omitted
wherever they were not confirmed. Entries should be checked before this
notebook is used as a source.

\end{warning}

\section{Introduction}
\label{sec:introduction}

Let $\mu$ denote the M\"obius function and let
\begin{equation}
M(x) = \sum_{n \le x} \mu(n)
\label{eq:mertens_definition}
\end{equation}
denote the Mertens function. We regard $M$ as a $\pm 1$ walk supported
on the squarefree integers, and we study it through three successive
conditionings: on the squarefree mask, on the number of sign changes,
and on the quadratic profile $n \mapsto M(n)^2$.

The principal object of the notebook is the third conditioning. We show
in Section~\ref{sec:reflection_fibre} that the fibre of the quadratic
profile is exactly an orbit of an elementary abelian reflection group
generated by excursion reflections, and that this group provides a
canonical parameter-free null model. We then carry a symmetric bilinear
form built from the divisor defect down to the quotient, obtaining a
form on the orientation bits of the excursion decomposition.

The motivation is a dichotomy that we prove in
Section~\ref{sec:blindness}: the quadratic profile is precisely the
invariant part of the two-point data under the reflection group, and
every cross-excursion two-point contribution to the quadratic energy
vanishes identically. Consequently no statistic that factors through
the quadratic profile can detect cross-excursion correlation of $\mu$.
The Gram form on orientation bits is introduced because it is, by
construction, reflection-odd, and therefore lives exactly in the
complement of what the energy can see.

Sections~\ref{sec:sine_kernel} and~\ref{sec:two_point} examine how far
this complement reaches toward, respectively, sine-kernel statistics
and two-point cancellation. Both sections end in conjectures, not
theorems. Section~\ref{sec:negative_results} records the negative
results and failed routes that shaped the construction, which we regard
as the most reliable content of the notebook.

Four ingredients are taken as inputs throughout, and are referred to by
the labels \textbf{(N1)}--\textbf{(N4)}:

\begin{enumerate}[label=\textbf{(N\arabic*)}]
\item a factorisation of the covariance of prefix squares for the
fixed-count sign bridge, yielding an $O(m)$ evaluation of the variance;
\item a single model combining a fixed squarefree mask, a fixed sign
count, and the M\"obius trajectory;
\item the formulation of the fibre of the quadratic profile as a
reflection orbit, together with a Gram form of the divisor defect on
the orientation bits;
\item an exact elementary form of a Robin-type equation for a truncated
functional.
\end{enumerate}

\textbf{(N3)} is the central ingredient and is developed in
Sections~\ref{sec:reflection_fibre}--\ref{sec:gram_form}.
\textbf{(N1)} and \textbf{(N2)} supply the null model and its error
bars in Sections~\ref{sec:three_conditionings}
and~\ref{sec:scale_resolved}. \textbf{(N4)} enters only in
Section~\ref{sec:robin_interface}, where it is treated as an interface
rather than as a tool.

\section{The three-constraint model}
\label{sec:three_conditionings}

\subsection{Mask, bits, and the fixed-count bridge}
\label{subsec:mask_and_bits}

\begin{definition}
\label{def:mask}
Fix $N \ge 2$. The \emph{squarefree mask} at level $N$ is
\begin{equation}
\mathcal{M}_N = \{\, n \le N : \mu(n) \neq 0 \,\},
\qquad
m = m(N) = |\mathcal{M}_N| .
\label{eq:mask_definition}
\end{equation}
The \emph{orientation bits} of the M\"obius path are the vector
\begin{equation}
\varepsilon = (\varepsilon_n)_{n \in \mathcal{M}_N} \in \{\pm 1\}^{m},
\qquad
\varepsilon_n = \mu(n).
\label{eq:orientation_bits}
\end{equation}
\end{definition}

With this notation $M(n) = \sum_{k \le n,\ k \in \mathcal{M}_N}
\varepsilon_k$ for $n \le N$, so that the mask carries all of the
multiplicative support information and $\varepsilon$ carries all of the
sign information.

\begin{definition}
\label{def:fixed_count_bridge}
Let $m_{\pm}$ denote the number of $n \in \mathcal{M}_N$ with
$\mu(n) = \pm 1$, so that $m_+ + m_- = m$ and $m_+ - m_- = M(N)$. The
\emph{fixed-count sign bridge} is the uniform probability measure
$\pi$ on the set of $\binom{m}{m_+}$ vectors in $\{\pm 1\}^m$ with
exactly $m_+$ entries equal to $+1$. We write
\begin{equation}
\bar{\varepsilon} = \frac{M(N)}{m}
\label{eq:mean_bit}
\end{equation}
for the common mean of the entries under $\pi$.
\end{definition}

Definition~\ref{def:fixed_count_bridge} is ingredient \textbf{(N2)}:
the mask is frozen, the sign count is frozen, and only the arrangement
of signs along the trajectory is randomised.

\subsection{Exact second moment of the bridge}
\label{subsec:bridge_second_moment}

\begin{proposition}
\label{prop:bridge_covariance}
Under the fixed-count sign bridge $\pi$ of
Definition~\ref{def:fixed_count_bridge},
\begin{equation}
\mathbb{E}_\pi\!\left[\varepsilon \varepsilon^{\mathsf{T}}\right]
= (1-\beta)\, I + \beta\, J,
\qquad
\beta = \frac{m \bar{\varepsilon}^{\,2} - 1}{m-1},
\label{eq:bridge_covariance}
\end{equation}
where $I$ is the identity and $J$ the all-ones matrix of size $m$.
\end{proposition}

\begin{proof}
Diagonal entries are $\varepsilon_n^2 = 1$ almost surely. For
$n \neq n'$, counting ordered pairs of distinct positions,
\begin{equation}
\mathbb{E}_\pi[\varepsilon_n \varepsilon_{n'}]
= \frac{m_+(m_+-1) + m_-(m_--1) - 2 m_+ m_-}{m(m-1)}
= \frac{(m_+-m_-)^2 - m}{m(m-1)} .
\label{eq:offdiagonal_computation}
\end{equation}
Since $m_+ - m_- = M(N) = m\bar{\varepsilon}$, the numerator equals
$m(m\bar{\varepsilon}^{\,2} - 1)$, which gives $\beta$ as stated.
\end{proof}

\begin{corollary}
\label{cor:bridge_projection_mean}
Let $P$ be an orthogonal projection on $\mathbb{R}^m$ of rank $d$ and
let $\mathbf{1}$ be the all-ones vector. Then
\begin{equation}
\mathbb{E}_\pi\!\left[\langle \varepsilon, P \varepsilon\rangle\right]
= (1-\beta)\, d + \beta\, \|P\mathbf{1}\|^2 .
\label{eq:bridge_projection_mean}
\end{equation}
If moreover $M(N) = o(\sqrt{m})$, then $\beta = O(1/m)$ and
\begin{equation}
\mathbb{E}_\pi\!\left[\langle \varepsilon, P \varepsilon\rangle\right]
= d + O(1).
\label{eq:bridge_projection_mean_asymptotic}
\end{equation}
\end{corollary}

\begin{proof}
Take traces against $P$ in \eqref{eq:bridge_covariance}, using
$\tr P = d$ and $\tr(PJ) = \|P\mathbf{1}\|^2$. For the second
statement, $|\beta| \le 1/(m-1)$ as soon as
$M(N)^2 \le m$, and $\|P\mathbf{1}\|^2 \le m$.
\end{proof}

Equation~\eqref{eq:bridge_projection_mean_asymptotic} is the structural
fact we shall use repeatedly: \emph{under the bridge, the expected
energy of an eigenspace is its dimension.} Nothing arithmetic enters.

\begin{remark}
\label{rem:finite_rank}
The right-hand side of \eqref{eq:bridge_covariance} is a scalar matrix
plus a rank-one correction. In particular the null second-moment
operator has finite rank above a multiple of the identity, and
therefore possesses no line spectrum in any auxiliary variable
conjugate to the trajectory. This is the structural content of
\textbf{(N1)} that the notebook uses; see
Remark~\ref{rem:N1_usage}.
\end{remark}

\begin{remark}
\label{rem:N1_usage}
Ingredient \textbf{(N1)}, the factorisation of
$\Cov_\pi\!\left(M(n)^2, M(n')^2\right)$ through the three-block
decomposition determined by $(n, n'-n, N-n')$ and the consequent $O(m)$
evaluation of $\Var_\pi$, is recorded in separate work and is not
reproduced here. This notebook uses only its structural consequence,
namely that the covariance kernel is separable of bounded rank, and
hence that the variance of any linear statistic of the prefix squares
is a finite combination of power sums of the mask counting function.
The explicit coefficients play no role below.
\end{remark}

\section{The reflection fibre of the quadratic profile}
\label{sec:reflection_fibre}

This section and the next develop ingredient \textbf{(N3)}.

\begin{definition}
\label{def:quadratic_profile}
The \emph{quadratic profile} of the path at level $N$ is the vector
\begin{equation}
\Pi_N(\varepsilon) = \left( M(n)^2 \right)_{n \le N}.
\label{eq:quadratic_profile}
\end{equation}
\end{definition}

\begin{definition}
\label{def:excursions}
Let
\begin{equation}
Z_N = \{\, n \le N : M(n) = 0 \,\} \cup \{0\}
\label{eq:zero_set}
\end{equation}
be the set of return points, written
$0 = z_0 < z_1 < \dots < z_{Z}$, where $Z = Z(N)$. The
\emph{excursions} are the blocks
\begin{equation}
E_j = \{\, n : z_{j-1} < n \le z_j \,\},
\qquad j = 1, \dots, Z,
\label{eq:excursion_blocks}
\end{equation}
together with the terminal block
$E_{Z+1} = \{\, n : z_{Z} < n \le N \,\}$, which is empty precisely
when $M(N) = 0$.
\end{definition}

\begin{definition}
\label{def:reflection_group}
The \emph{excursion reflection group} is
\begin{equation}
\mathcal{R}_N \;\cong\; (\mathbb{Z}/2)^{\,Z(N)},
\label{eq:reflection_group}
\end{equation}
acting on orientation bits by
\begin{equation}
(\sigma \cdot \varepsilon)_n = \sigma_{j(n)}\, \varepsilon_n,
\qquad \sigma = (\sigma_1,\dots,\sigma_Z) \in \{\pm 1\}^{Z},
\label{eq:reflection_action}
\end{equation}
where $j(n)$ is the index of the excursion containing $n$, and where
the terminal block is held fixed.
\end{definition}

\begin{proposition}
\label{prop:reflection_fibre}
The action \eqref{eq:reflection_action} preserves the quadratic
profile. More precisely, for every
$\sigma \in \mathcal{R}_N$ and every $n \le N$,
\begin{equation}
M_{\sigma \cdot \varepsilon}(n) = \sigma_{j(n)}\, M_\varepsilon(n),
\label{eq:path_equivariance}
\end{equation}
and consequently
$\Pi_N(\sigma\cdot\varepsilon) = \Pi_N(\varepsilon)$. Conversely, two
sign vectors supported on the same mask with the same quadratic profile
differ by an element of $\mathcal{R}_N$. Hence the fibre of $\Pi_N$
through $\varepsilon$ is exactly the orbit
$\mathcal{R}_N \cdot \varepsilon$.
\end{proposition}

\begin{proof}
For $n \in E_j$ we have $M_\varepsilon(z_{j-1}) = 0$, so
$M_\varepsilon(n) = \sum_{z_{j-1} < k \le n} \varepsilon_k$, a sum
entirely inside $E_j$; applying $\sigma$ multiplies every summand by
$\sigma_j$, which gives \eqref{eq:path_equivariance} and hence
invariance of $M(n)^2$.

Conversely, suppose $\varepsilon$ and $\varepsilon'$ have the same mask
and satisfy $M_{\varepsilon}(n)^2 = M_{\varepsilon'}(n)^2$ for all
$n \le N$. Then the two paths have the same zero set, so the same
excursion decomposition, and $|M_\varepsilon| = |M_{\varepsilon'}|$
pointwise. On each excursion $E_j$ the path does not vanish in the
interior, so $M_\varepsilon$ has constant sign there; define $\sigma_j$
to be the ratio of the two signs. Since a path on a fixed mask is
determined by its increments, and the increments of
$M_{\varepsilon'}$ on $E_j$ are $\sigma_j$ times those of
$M_\varepsilon$, we obtain $\varepsilon' = \sigma \cdot \varepsilon$.
\end{proof}

\begin{remark}
\label{rem:catalan_link}
Proposition~\ref{prop:reflection_fibre} identifies the combinatorial
object underlying the excursion (Dyck) decomposition of the M\"obius
walk with a reflection group in the classical sense. The reflection
principle for lattice paths is standard; what is used here is not the
principle itself but the observation that the group it generates is
exactly the isotropy group of the quadratic profile.
\end{remark}

\subsection{The canonical null on orientation bits}
\label{subsec:canonical_null}

\begin{definition}
\label{def:excursion_signs}
Given $\varepsilon$, its \emph{excursion signs} are
$s = (s_1,\dots,s_Z) \in \{\pm 1\}^Z$, where $s_j$ records the sign of
$M_\varepsilon$ in the interior of $E_j$. The \emph{reflection null}
$\nu$ is the uniform probability measure on $\{\pm 1\}^Z$.
\end{definition}

\begin{proposition}
\label{prop:exact_isotropy}
Under the reflection null $\nu$ one has
\begin{equation}
\mathbb{E}_\nu\!\left[s s^{\mathsf{T}}\right] = I_Z
\label{eq:exact_isotropy}
\end{equation}
exactly, with no rank-one correction. Consequently, for every symmetric
$A$ on $\mathbb{R}^Z$,
\begin{equation}
\mathbb{E}_\nu\!\left[\langle s, A s\rangle\right] = \tr A,
\qquad
\Var_\nu\!\left[\langle s, A s\rangle\right]
= 2 \sum_{j \neq k} A_{jk}^2 .
\label{eq:exact_mean_variance}
\end{equation}
In particular, if $A = \widehat{P}$ is an orthogonal projection of rank
$d$, then
\begin{equation}
\mathbb{E}_\nu\!\left[\langle s, \widehat{P} s\rangle\right] = d,
\qquad
\Var_\nu\!\left[\langle s, \widehat{P} s\rangle\right]
= 2\left( d - \sum_{j} \widehat{P}_{jj}^2 \right).
\label{eq:exact_projection_moments}
\end{equation}
\end{proposition}

\begin{proof}
The coordinates $s_j$ are independent Rademacher variables under
$\nu$, whence \eqref{eq:exact_isotropy}. Writing
$\langle s, As\rangle = \tr A + \sum_{j \neq k} A_{jk} s_j s_k$ and
using $\mathbb{E}_\nu[s_j s_k s_{j'} s_{k'}] = 1$ exactly when
$\{j,k\} = \{j',k'\}$ and $0$ otherwise gives the variance. For a
projection, $\sum_{j,k}\widehat{P}_{jk}^2 = \tr \widehat{P}^2
= \tr \widehat{P} = d$, which yields
\eqref{eq:exact_projection_moments}.
\end{proof}

\begin{remark}
\label{rem:hierarchy_of_nulls}
Corollary~\ref{cor:bridge_projection_mean} and
Proposition~\ref{prop:exact_isotropy} form a hierarchy. Conditioning on
the mask removes the squarefree density; conditioning on the sign count
removes the endpoint drift; conditioning in addition on the quadratic
profile leaves exactly the reflection group, whose null is isotropic
\emph{without} the $O(1)$ correction of
\eqref{eq:bridge_projection_mean_asymptotic}. The third conditioning is
therefore not a refinement of convenience: it is the step at which the
null becomes parameter-free.
\end{remark}

\section{The divisor-defect Gram form on orientation bits}
\label{sec:gram_form}

We now introduce the second half of ingredient \textbf{(N3)}. Because
the notebook uses only structural properties of the form, we isolate
those properties as hypotheses rather than fixing a normalisation.

\begin{definition}
\label{def:divisor_defect_form}
Let $B$ denote the divisor-defect matrix attached to the mask
$\mathcal{M}_N$, and let
\begin{equation}
G = B^{\mathsf{T}} B
\label{eq:gram_form}
\end{equation}
be its Gram form, acting on $\mathbb{R}^{m}$ indexed by the mask.
\end{definition}

\begin{definition}
\label{def:hypotheses_on_G}
We say that $G$ is \emph{admissible} if it satisfies:
\begin{enumerate}[label=(G\arabic*)]
\item $G$ is real symmetric and positive semidefinite;
\item the eigenbasis of $G$ factorises over primes, in the sense that
the associated kernel admits a tensor factorisation
\begin{equation}
\widehat{K}(r,q) = \prod_{p} L_p\!\left( v_p(r), v_p(q) \right)
\label{eq:tensor_factorisation}
\end{equation}
into local factors indexed by the prime valuations.
\end{enumerate}
\end{definition}

Property (G2) is an input from separate work and is not proved here. It
is the only place where the arithmetic of the divisor defect enters the
constructions below, and its role is precisely to guarantee that the
eigenbasis of $G$ is a character basis of the multiplicative structure;
see Section~\ref{sec:sine_kernel}.

\begin{definition}
\label{def:pushed_form}
The \emph{orientation-bit form} is the $Z \times Z$ symmetric matrix
\begin{equation}
\widehat{G}_{jk} \;=\; \sum_{a \in E_j} \sum_{b \in E_k} G_{ab},
\qquad 1 \le j,k \le Z,
\label{eq:pushed_form}
\end{equation}
obtained by summing $G$ over excursion blocks. Its spectral resolution
is written
\begin{equation}
\widehat{G} = \sum_{k} \lambda_k \widehat{P}_k,
\qquad d_k = \rank \widehat{P}_k .
\label{eq:pushed_spectral}
\end{equation}
\end{definition}

\begin{proposition}
\label{prop:pushed_form_odd}
The quadratic form $s \mapsto \langle s, \widehat{G} s\rangle$ is
invariant under the global sign flip $s \mapsto -s$, and its
off-diagonal part is odd under every individual excursion reflection
$s_j \mapsto -s_j$. In particular
$\langle s, \widehat{G} s\rangle$ is not a function of the quadratic
profile unless $\widehat{G}$ is diagonal.
\end{proposition}

\begin{proof}
Immediate from \eqref{eq:pushed_form}: the term $s_j s_k$ changes sign
under $s_j \mapsto -s_j$ whenever $j \neq k$, while the diagonal terms
$\widehat{G}_{jj} s_j^2 = \widehat{G}_{jj}$ are constant. By
Proposition~\ref{prop:reflection_fibre} the quadratic profile is
constant along $\mathcal{R}_N$-orbits, so any function of the profile
is $\mathcal{R}_N$-invariant.
\end{proof}

Proposition~\ref{prop:pushed_form_odd} is the reason the form
\eqref{eq:pushed_form} is introduced. The next section shows that this
is not merely a convenient property but the only available one.

\section{Blindness of the quadratic profile}
\label{sec:blindness}

\begin{lemma}
\label{lem:energy_as_two_point}
For every $N \ge 1$,
\begin{equation}
\sum_{n \le N} M(n)^2
= \sum_{a,b \le N} \mu(a)\mu(b) \left( N + 1 - \max(a,b) \right).
\label{eq:energy_as_two_point}
\end{equation}
\end{lemma}

\begin{proof}
Expand $M(n)^2 = \sum_{a,b \le n} \mu(a)\mu(b)$ and exchange the order
of summation; the pair $(a,b)$ occurs for every $n$ with
$\max(a,b) \le n \le N$, that is for exactly $N + 1 - \max(a,b)$ values
of $n$.
\end{proof}

Lemma~\ref{lem:energy_as_two_point} exhibits the quadratic energy as a
triangularly weighted aggregate of two-point data. The following result
shows that this aggregate is blind to every cross-excursion pair.

\begin{theorem}
\label{thm:cross_excursion_vanishing}
Assume $M(N) = 0$, so that the excursion decomposition of
Definition~\ref{def:excursions} has no terminal block. Then for all
$j \neq k$,
\begin{equation}
\sum_{a \in E_j} \sum_{b \in E_k}
\mu(a)\mu(b)\left( N + 1 - \max(a,b) \right) = 0 .
\label{eq:cross_excursion_vanishing}
\end{equation}
Consequently the quadratic energy $\sum_{n \le N} M(n)^2$ depends only
on the within-excursion two-point data
$\{\mu(a)\mu(b) : a, b \in E_j,\ 1 \le j \le Z\}$.
\end{theorem}

\begin{proof}
Let $j < k$, so that every $a \in E_j$ precedes every $b \in E_k$ and
hence $\max(a,b) = b$. The double sum factorises as
\begin{equation}
\left( \sum_{a \in E_j} \mu(a) \right)
\left( \sum_{b \in E_k} \mu(b)\left(N+1-b\right) \right).
\label{eq:cross_excursion_factorisation}
\end{equation}
By Definition~\ref{def:excursions} the path returns to zero at both
endpoints of $E_j$, so $\sum_{a \in E_j} \mu(a) =
M(z_j) - M(z_{j-1}) = 0$, and the first factor vanishes. The case
$j > k$ is symmetric.
\end{proof}

\begin{corollary}
\label{cor:profile_blindness}
No statistic that factors through the quadratic profile $\Pi_N$ can
detect cross-excursion correlation of $\mu$. In particular the energy
$\sum_{n \le N} M(n)^2$, all of its moments, and all functionals of the
excursion heights and lengths are blind to it.
\end{corollary}

\begin{proof}
By Proposition~\ref{prop:reflection_fibre} a statistic factoring
through $\Pi_N$ is $\mathcal{R}_N$-invariant, hence constant on the
fibre; by Theorem~\ref{thm:cross_excursion_vanishing} the
cross-excursion two-point mass makes no contribution to the energy, and
by \eqref{eq:reflection_action} it is precisely the part of the
two-point data on which $\mathcal{R}_N$ acts nontrivially.
\end{proof}

\begin{remark}
\label{rem:blindness_interpretation}
Theorem~\ref{thm:cross_excursion_vanishing} is elementary but it
delimits an entire family of approaches. Energy-type and
excursion-geometry approaches to the M\"obius walk are confined, by the
structure of the return points, to within-excursion information. The
vanishing in \eqref{eq:cross_excursion_vanishing} is caused by the
defining property of an excursion, namely that its increments sum to
zero, so it cannot be repaired by reweighting: any weight depending on
$a$ and $b$ only through $\max(a,b)$ produces the same factorisation
\eqref{eq:cross_excursion_factorisation}.
\end{remark}

\begin{remark}
\label{rem:terminal_block}
If $M(N) \neq 0$ the terminal block $E_{Z+1}$ is nonempty and
\eqref{eq:cross_excursion_vanishing} fails for pairs involving it,
since $\sum_{a \in E_{Z+1}} \mu(a) = M(N)$. The resulting correction is
a single rank-one term proportional to $M(N)$. Throughout the notebook
we normalise by taking $N$ at a return point; the general case differs
by that term.
\end{remark}

\section{The scale-resolved statistic}
\label{sec:scale_resolved}

\begin{definition}
\label{def:scale_resolved_statistic}
With $\widehat{G} = \sum_k \lambda_k \widehat{P}_k$ as in
\eqref{eq:pushed_spectral} and $s$ the excursion signs of the true
M\"obius path, set
\begin{equation}
\mathcal{E}_k = \langle s, \widehat{P}_k s \rangle,
\qquad
\mathcal{R}_k = \frac{\mathcal{E}_k}{d_k}.
\label{eq:scale_resolved_statistic}
\end{equation}
\end{definition}

\begin{proposition}
\label{prop:statistic_calibration}
Under the reflection null $\nu$,
\begin{equation}
\mathbb{E}_\nu\!\left[\mathcal{R}_k\right] = 1
\quad\text{exactly, and}\quad
\Var_\nu\!\left[\mathcal{R}_k\right]
= \frac{2}{d_k^2}\left( d_k - \sum_j \widehat{P}_{k,jj}^2 \right).
\label{eq:statistic_calibration}
\end{equation}
Moreover $\sum_k \mathcal{E}_k = \|s\|^2 = Z$ and $\sum_k d_k = Z$, so
the statistic is a partition of a conserved total.
\end{proposition}

\begin{proof}
Immediate from Proposition~\ref{prop:exact_isotropy} applied to
$A = \widehat{P}_k$, together with $\sum_k \widehat{P}_k = I_Z$.
\end{proof}

\begin{remark}
\label{rem:no_free_parameters}
The calibration \eqref{eq:statistic_calibration} contains no free
parameter, no fitted constant and no asymptotic normalisation. This is
a direct consequence of the third conditioning; compare
Remark~\ref{rem:hierarchy_of_nulls}. By
Proposition~\ref{prop:pushed_form_odd} and
Corollary~\ref{cor:profile_blindness}, $\mathcal{R}_k$ is a
reflection-odd statistic and therefore measures exactly the information
that the quadratic profile cannot carry.
\end{remark}

\begin{definition}
\label{def:anisotropy}
The \emph{Gram anisotropy} of the M\"obius path at level $N$ is the
profile $k \mapsto \mathcal{R}_k - 1$.
\end{definition}

\begin{remark}
\label{rem:empirical_inputs}
The following numerical quantities are recorded in separate work and
are stated here only to indicate the scale of the effects that a future
analysis of the anisotropy profile would have to account for. They are
not computed in this notebook and no argument below depends on their
values: a stable deficit of roughly $0.137$ in the ratio of the true
quadratic energy to its permutation-null mean; a defect of
approximately $12\%$ at order two and $32\%$ at order three in the test
of asymptotic freeness of the arithmetic and shift operators; a
persistent bias of about $2.5\%$ in a decomposition with a
parameter-free arithmetic factor; and the systematic observation that
the anomaly at $d=1$ is the weakest, with stronger decorrelation for
$d \ge 2$.
\end{remark}

\begin{conjecture}
\label{conj:anisotropy_nontrivial}
The Gram anisotropy is not identically $o(1)$: there exist eigenspaces
of $\widehat{G}$ on which
$\mathcal{R}_k - 1$ exceeds the null standard deviation
\eqref{eq:statistic_calibration} by an unbounded factor as
$N \to \infty$.
\end{conjecture}

\begin{heuristic}
\label{heur:defects_are_one_object}
The four empirical quantities listed in
Remark~\ref{rem:empirical_inputs} are heuristically interpretable as
four measurements of a single quantity, namely the deviation of
$\widehat{G}$ from a flat band projector. This is a heuristic
identification; it is offered because it is falsifiable, in the
following sense: it predicts that the $d$-dependence of the anisotropy
profile reproduces the recorded pattern in which $d=1$ is the weakest
anomaly. If the anisotropy profile does not exhibit that pattern, the
identification is false.
\end{heuristic}

\section{Route toward sine-kernel statistics}
\label{sec:sine_kernel}

\subsection{Flat bands and the sine kernel}
\label{subsec:flat_bands}

We record the standard fact that isolates what a derivation of
sine-kernel statistics must supply \cite{Mehta2004}.

\begin{lemma}
\label{lem:band_projector_sinc}
Let $P$ be the orthogonal projection onto a contiguous block of $d$
characters of a cyclic group of order $Z$. In the position basis the
kernel of $P$ is a Dirichlet kernel, and under the scaling
$Z \to \infty$ with $d/Z$ fixed and positions rescaled by the mean
spacing, it converges to the sine kernel
$\sin(\pi x)/(\pi x)$.
\end{lemma}

\begin{proof}
The kernel is $P(n,n') = \frac{1}{Z}\sum_{t \in S}
e^{2\pi i t (n-n')/Z}$ with $S$ an interval of length $d$, which is a
Dirichlet kernel; the stated limit is the standard scaling limit of the
Dirichlet kernel. See \cite{Mehta2004} for the random matrix context.
\end{proof}

The content of Lemma~\ref{lem:band_projector_sinc} is that the sine
kernel is the kernel of a \emph{flat band projector}. Equivalently, on
the Fourier side, the unitary form factor $\min(|\alpha|,1)$ is the
indicator of a band, and the pair correlation of Montgomery
\cite{Montgomery1973} is known in the range $|\alpha| < 1$ while the
plateau $|\alpha| > 1$ is open.

\begin{warning}
\label{warn:sinc_is_the_null}
By Proposition~\ref{prop:exact_isotropy} the reflection null has
$\mathbb{E}_\nu[ss^{\mathsf{T}}] = I$, which restricted to any
contiguous block of characters is, by
Lemma~\ref{lem:band_projector_sinc}, a sine kernel. Hence
\emph{the sine kernel is the null hypothesis of the statistic
$\mathcal{R}_k$, not a consequence of it.} A derivation of sine-kernel
behaviour from the flatness of the null would be a derivation of the
null from the null, and carries no evidential weight concerning the
zeros of $\zeta$. Any claim that this framework "produces the sine
kernel" must be read against this warning.
\end{warning}

\begin{remark}
\label{rem:multiplicative_band}
Hypothesis (G2) of Definition~\ref{def:hypotheses_on_G} states that the
eigenbasis of $G$ factorises over primes, hence is a character basis of
the multiplicative structure rather than of an additive cyclic group.
Consequently the band of Lemma~\ref{lem:band_projector_sinc}, when
transported to positions, produces a sine kernel in the variable
$\log n$ rather than in $n$. This is the variable in which the explicit
formula oscillates, since $n^{i\gamma} = e^{i\gamma \log n}$, so the
identification is the natural one; it does not, however, affect
Warning~\ref{warn:sinc_is_the_null}.
\end{remark}

\subsection{The symmetry obstruction and the twist}
\label{subsec:symmetry_obstruction}

\begin{remark}
\label{rem:symmetry_class_obstruction}
The reflection group $\mathcal{R}_N \cong (\mathbb{Z}/2)^Z$ and the
form $\widehat{G}$ are real. A real symmetric ensemble has
$\beta = 1$ (orthogonal) statistics, whose form factor
$b(\tau) = 2\tau - \tau\log(1+2\tau)$ on $\tau \le 1$
\cite{Mehta2004} is not the indicator of a band and has no plateau;
the orthogonal two-point function contains the sine kernel but is not
$1 - \left(\sin \pi x / \pi x\right)^2$. The conjectural symmetry class
for the zeros of $\zeta$ is $\beta = 2$
\cite{Montgomery1973, Odlyzko1987}. Therefore the present construction,
taken as it stands, is in the wrong symmetry class, and the obstruction
is structural rather than quantitative.
\end{remark}

The natural way to leave the orthogonal class inside the present
framework is to complexify the reflection group. We state this as a
conjecture.

\begin{conjecture}
\label{conj:twisted_orientation_bits}
Let $U(1)^Z$ act on orientation bits by
$s_j \mapsto e^{i\vartheta_j} s_j$, extending
\eqref{eq:reflection_action} from
$(\mathbb{Z}/2)^Z \subset U(1)^Z$, and let
$\widehat{G}^{(\vartheta)}$ denote the Hermitian form obtained by
transporting $G$ through the twisted action. Then as $\vartheta$ varies
the local statistics of $\widehat{G}^{(\vartheta)}$ interpolate between
the orthogonal and unitary classes, and the crossover occurs at a
critical scale $\vartheta_c \asymp 1/Z$, namely when the
antisymmetric part of the form is comparable to the mean level spacing.
\end{conjecture}

\begin{remark}
\label{rem:pandey_mehta}
Conjecture~\ref{conj:twisted_orientation_bits} is modelled on the
known interpolating ensembles between the orthogonal and unitary
classes \cite{PandeyMehta1983}; what is conjectural here is not the
existence of such an interpolation in random matrix theory but the
claim that the twist of the orientation bits realises it for
$\widehat{G}$, and that the critical scale is as stated.
\end{remark}

\begin{warning}
\label{warn:crossover_is_necessary_not_sufficient}
Even if Conjecture~\ref{conj:twisted_orientation_bits} holds, it would
establish only that the ensemble is in a symmetry class compatible with
the zeros of $\zeta$. Compatibility of symmetry class is a necessary
condition and is very far from an implication about
$\zeta$.
\end{warning}

\begin{question}
\label{q:plateau_as_flatness}
Is there a formulation of the plateau $F(\alpha) = 1$ for $\alpha > 1$
as a flatness statement for $\widehat{G}^{(\vartheta)}$ on a band of
eigenspaces? A positive answer would be a reformulation rather than a
proof, but it would move the plateau to a setting carrying a transfer
structure along the excursion index, which the arithmetic side does not
provide.
\end{question}

\section{Route toward two-point cancellation}
\label{sec:two_point}

\begin{remark}
\label{rem:chowla_lives_in_the_odd_part}
Combining Lemma~\ref{lem:energy_as_two_point},
Theorem~\ref{thm:cross_excursion_vanishing} and
Proposition~\ref{prop:pushed_form_odd}, the two-point data of $\mu$
splits into a reflection-even part, carried entirely by the quadratic
profile and entirely within excursions, and a reflection-odd part,
consisting of all cross-excursion pairs. Two-point cancellation for
$\mu$ in the sense of Chowla's conjecture is a statement about shifted
pairs $\mu(a)\mu(a+h)$ with $h$ unrestricted, and such pairs are
cross-excursion for all but a vanishing proportion of $a$ once $h$
exceeds typical excursion lengths. Hence the content of two-point
cancellation lies in the reflection-odd part, which is precisely the
part measured by $\mathcal{R}_k$.
\end{remark}

\begin{conjecture}
\label{conj:cross_excursion_mass}
Let
\begin{equation}
C_N(h) = \sum_{a \le N - h} \mu(a)\mu(a+h)
\label{eq:shifted_correlation}
\end{equation}
denote the shifted correlation. Then the aggregate cross-excursion mass
of $\{C_N(h)\}_h$ is, in the mean, of the order predicted by the
reflection null $\nu$; equivalently, the Gram anisotropy of
Definition~\ref{def:anisotropy} is bounded in the relevant eigenspaces.
\end{conjecture}

\begin{heuristic}
\label{heur:logarithmic_weighting}
The triangular weight $N+1-\max(a,b)$ of
\eqref{eq:energy_as_two_point} aggregates all shifts $h$ up to $N$ with
comparable mass. Logarithmic weighting is the regime in which two-point
cancellation for multiplicative functions is presently accessible
\cite{Tao2016}, building on the short-interval theory of
\cite{MatomakiRadziwill2016}. This suggests, heuristically, replacing
the triangular weight by a logarithmic one in
\eqref{eq:energy_as_two_point} before pushing to orientation bits. The
obstruction that then remains is uniformity in $h$: the available
theorems are for fixed or slowly growing $h$, whereas the aggregate
requires $h$ up to $N$. We do not claim that this gap can be closed;
we record that it is, on this route, the only gap.
\end{heuristic}

\begin{question}
\label{q:odd_statistic_and_uniformity}
Does the passage to the reflection-odd statistic $\mathcal{R}_k$ change
the uniformity requirement in
Heuristic~\ref{heur:logarithmic_weighting}? Since
$\widehat{G}$ aggregates pairs by excursion block rather than by shift,
the relevant uniformity is in the block index rather than in $h$, and
the two need not be equivalent.
\end{question}

\section{The Robin interface}
\label{sec:robin_interface}

Ingredient \textbf{(N4)} enters here. The form $G$ of
Definition~\ref{def:divisor_defect_form} is built from the divisor
defect; Robin's criterion \cite{Robin1984} is an inequality on the
divisor sum $\sigma(n)$, equivalent to the Riemann hypothesis, and
admits elementary reformulations \cite{Lagarias2002}. The two therefore
concern the same arithmetic object from opposite directions.

\begin{question}
\label{q:robin_cone_comparison}
Suppose the truncated functional attached to $G$ satisfies an exact
elementary Robin-type equation. Does the positivity cone of that
equation coincide with the cone associated with the tail-stability
criterion of the ordinal-prime deformation? If the cones coincide, the
two criteria are reparametrisations of one another; if they differ, the
pair yields a positivity family not reducible to either.
\end{question}

\begin{warning}
\label{warn:robin_novelty}
The literature on Robin-type criteria is extensive. No claim of novelty
is made for \textbf{(N4)} in this notebook, and the exact elementary
form should be checked against the existing literature before being
described as new. Question~\ref{q:robin_cone_comparison} is stated as an
interface, not as a result.
\end{warning}

\section{Negative results and failed routes}
\label{sec:negative_results}

In accordance with the archival principle we record the routes that
were examined and abandoned, since the reasons are more durable than
the attempts.

\begin{remark}
\label{rem:negative_result_energy_scalar}
The scalar quadratic energy, evaluated at a single level $N$, carries
no local information. Its cross terms appear with a weight that is
essentially constant on the scale of the mean spacing of zero
ordinates, so no single value of the energy can resolve that scale. The
defect is not in the weight but in the collapse to a single number: as
a function of $\log N$ the energy is almost periodic with frequency set
given by differences of ordinates, and the associated multiplier is
invertible at every real frequency. This route is therefore not closed,
but it is subsumed by the present framework, in which the same
information appears without the auxiliary variable.
\end{remark}

\begin{remark}
\label{rem:negative_result_height_weighting}
Within the family of statistics built from $M$ alone, the weights
attached to individual zeros decay in the ordinate, so the resulting
pair-correlation measure mixes heights and is not a local statistic.
This cannot be repaired by any operation on $\log N$, since that
variable is conjugate to differences of ordinates and not to the
ordinates themselves. The present notebook avoids the difficulty by not
working in ordinate variables at all.
\end{remark}

\begin{remark}
\label{rem:negative_result_jacobi}
The excursion decomposition admits a classical passage to continued
fractions and hence to a Jacobi matrix \cite{Flajolet1980}. That route
produces a real symmetric operator and therefore lands in the
orthogonal symmetry class, for the same structural reason discussed in
Remark~\ref{rem:symmetry_class_obstruction}. It is recorded because the
coincidence of the two obstructions, in the excursion picture and in the
Gram picture, is what motivates
Conjecture~\ref{conj:twisted_orientation_bits}.
\end{remark}

\begin{remark}
\label{rem:negative_result_topology}
A cell-complex description of the Mertens function exists, with
$M(n)$ recovered as an Euler characteristic and Betti numbers given by
odd divisor sums \cite{Bjorner2011}. It does not reach two-point
cancellation, because the complex is defined multiplicatively while
shifted correlations are additive: there is no map relating the complex
at $n$ to the complex at $n+h$. Constructing a complex indexed by pairs
$(n, n+h)$ whose Euler characteristic is $C_N(h)$ would remove the
obstruction; we are not aware of such an object, and we make no claim
that it exists.
\end{remark}

\begin{remark}
\label{rem:negative_result_omega_lower_bound}
Any hierarchy of bounds on the M\"obius walk must respect the
unconditional oscillation result $M(x) = \Omega_\pm(\sqrt{x})$ with
explicit constants \cite{OdlyzkoteRiele1985}. In particular no
$\varepsilon$ appearing in such a hierarchy is removable, and a
framework that appears to remove one has an error.
\end{remark}

\section{Open questions and future work}
\label{sec:open_questions}

\begin{enumerate}[label=(Q\arabic*)]
\item Establish or refute Conjecture~\ref{conj:anisotropy_nontrivial},
and if it holds, determine which eigenspaces of $\widehat{G}$ carry the
anisotropy.
\item Test Heuristic~\ref{heur:defects_are_one_object} through the
predicted $d$-dependence of the anisotropy profile. This is the
cheapest falsification available in the framework.
\item Decide whether property (G2) of
Definition~\ref{def:hypotheses_on_G} survives the passage to
orientation bits, that is, whether $\widehat{G}$ inherits a band
structure from $G$ or whether the excursion aggregation destroys it.
This is the pivotal technical question of the notebook and is
completely open.
\item Determine the growth of $Z(N)$, the number of return points. Every
statement in Sections~\ref{sec:scale_resolved}--\ref{sec:two_point} is
uniform in $Z$, but the strength of the conclusions depends on it, and
we make no claim about it here.
\item Settle Question~\ref{q:plateau_as_flatness} and
Question~\ref{q:odd_statistic_and_uniformity}.
\item Settle Question~\ref{q:robin_cone_comparison}, subject to
Warning~\ref{warn:robin_novelty}.
\end{enumerate}

\section*{Notebook History}
\label{sec:notebook_history}

Version 1.0:
Initial version. Three-constraint model and exact bridge second moment;
reflection fibre of the quadratic profile and the canonical
parameter-free null; divisor-defect Gram form pushed to orientation
bits; blindness theorem for the quadratic profile; scale-resolved
statistic with exact mean and variance; exploratory routes toward
sine-kernel statistics and two-point cancellation; negative results
recorded.

\begin{thebibliography}{99}

\bibitem{Bjorner2011}
A.~Bj\"orner,
\emph{A cell complex in number theory},
arXiv:1101.5704.

\bibitem{Flajolet1980}
P.~Flajolet,
\emph{Combinatorial aspects of continued fractions},
Discrete Mathematics \textbf{32} (1980), 125--161.

\bibitem{Gonek1989}
S.~M.~Gonek,
\emph{On negative moments of the Riemann zeta-function},
Mathematika \textbf{36} (1989), 71--88.

\bibitem{Lagarias2002}
J.~C.~Lagarias,
\emph{An elementary problem equivalent to the Riemann hypothesis},
American Mathematical Monthly \textbf{109} (2002), 534--543.

\bibitem{MatomakiRadziwill2016}
K.~Matom\"aki and M.~Radziwi\l{}\l{},
\emph{Multiplicative functions in short intervals},
Annals of Mathematics \textbf{183} (2016), 1015--1056.

\bibitem{Mehta2004}
M.~L.~Mehta,
\emph{Random Matrices}, third edition,
Elsevier, Amsterdam, 2004.

\bibitem{Montgomery1973}
H.~L.~Montgomery,
\emph{The pair correlation of zeros of the zeta function},
in: Analytic Number Theory,
Proceedings of Symposia in Pure Mathematics \textbf{24},
American Mathematical Society, 1973, 181--193.

\bibitem{Odlyzko1987}
A.~M.~Odlyzko,
\emph{On the distribution of spacings between zeros of the zeta
function},
Mathematics of Computation \textbf{48} (1987), 273--308.

\bibitem{OdlyzkoteRiele1985}
A.~M.~Odlyzko and H.~J.~J.~te Riele,
\emph{Disproof of the Mertens conjecture},
Journal f\"ur die reine und angewandte Mathematik \textbf{357} (1985),
138--160.

\bibitem{PandeyMehta1983}
A.~Pandey and M.~L.~Mehta,
\emph{Gaussian ensembles of random Hermitian matrices intermediate
between orthogonal and unitary ones},
Communications in Mathematical Physics \textbf{87} (1983), 449--468.

\bibitem{Robin1984}
G.~Robin,
\emph{Grandes valeurs de la fonction somme des diviseurs et hypoth\`ese
de Riemann},
Journal de Math\'ematiques Pures et Appliqu\'ees \textbf{63} (1984),
187--213.

\bibitem{Tao2016}
T.~Tao,
\emph{The logarithmically averaged Chowla and Elliott conjectures for
two-point correlations},
Forum of Mathematics, Pi \textbf{4} (2016), e8.

\bibitem{Titchmarsh1986}
E.~C.~Titchmarsh,
\emph{The Theory of the Riemann Zeta-Function}, second edition,
revised by D.~R.~Heath-Brown,
Oxford University Press, 1986.

\end{thebibliography}

\label{sec:bibliography}

\end{document}
